\sect{Формальная постановка задачи}{problem-formal}
%
Прежде чем перейти к рассмотрению алгоритмов, договоримся об
обозначениях. Поскольку математический аппарат отслеживания
партитуры заимствован из четырёх соседних областей~--- теории
вероятностей, обработки сигналов, оптимального выравнивания и
обучения с подкреплением,~--- в литературе по нему используется
весьма пёстрый набор греческих и латинских букв. Чтобы избежать
путаницы при чтении, мы заранее перечисляем все ключевые
обозначения, по возможности повторяющие нотацию первоисточников.

\paragraph*{Список основных обозначений.}
\textit{Партитура и исполнение.} $S$~--- партитура (нотный
текст), $N$~--- число нот в ней; $i$ нумерует ноту партитуры от
$1$ до $N$. $P$~--- исполнение (живая запись), $M$~--- число
наблюдаемых событий (заранее неизвестно); $j$ нумерует событие
исполнения. $p_i$~--- высота $i$-й ноты партитуры в стандарте
MIDI, $q_j$~--- высота $j$-го события исполнения. Моменты
времени: $t_i$~--- номинальный (по партитуре), $\tau_j$~---
реальный (по исполнению). $v_0$~--- номинальный темп
партитуры (beat/секунду); $\hat v(t)$~--- текущая оценка темпа
исполнения.
\textit{Функция выравнивания.} Греческая $\varphi$ (читается
«фи») традиционно обозначает функцию выравнивания между
исполнением и партитурой~\cite{Cont2010,Mueller2007}; $\varphi(j)$~---
истинная позиция исполнителя на момент $j$-го события,
$\hat\varphi(j)$~--- её оценка (с «крышкой» по стандартной
статистической традиции).
\textit{Скрытая марковская модель} (по Рабинеру~\cite{Rabiner1989}):
$\lambda = (A, B, \pi)$, где $A=(A_{ij})$~--- матрица
переходов между состояниями, $B = \{b_i(o)\}$~--- эмиссионные
распределения, $\pi=(\pi_i)$~--- начальное распределение по
состояниям. $\alpha_t(j)$~--- прямая (forward) переменная
алгоритма, $\beta_t(j)$~--- обратная (backward) переменная
(не путать с темпом, для которого мы используем $v$);
$\delta_t(j)$, $\psi_t(j)$~--- переменные Viterbi-алгоритма;
$d_j(u)$~--- распределение длительности состояния в
полу-марковской модели~\cite{Cont2010}.
\textit{Обучение с подкреплением}: $\mathcal{M} = (\mathcal{S},
\mathcal{A}, \mathcal{T}, r, \gamma)$~--- марковский процесс
принятия решений; $s_t \in \mathcal{S}$~--- состояние,
$a_t \in \mathcal{A}$~--- действие, $r_t$~--- мгновенное
вознаграждение, $\gamma \in [0, 1]$~--- дисконт-фактор;
$\pi_\theta(a\mid s)$~--- параметризованная политика
(не путать с начальным распределением HMM).

С этой нотацией формализация задачи ниже остаётся короткой.

\paragraph*{Партитура.} Партитура задаётся последовательностью
$S = \{(p_i, t_i^{\text{on}}, t_i^{\text{off}})\}_{i=1}^{N}$ из
$N$~нот, в которой $p_i \in \{0, \dots, 127\}$~--- высота в
стандарте MIDI, а $t_i^{\text{on}}, t_i^{\text{off}}$~---
номинальные моменты начала и конца звучания относительно
базового темпа $v_0$ (как правило, в beat-единицах). Исполнение
поступает каузально как поток
$P = \{(q_j, \tau_j^{\text{on}}, \tau_j^{\text{off}})\}_{j=1}^{M}$
с реализованными высотами $q_j$ и реальными моментами времени
$\tau_j$.

\paragraph*{Задача.} \textit{Задача отслеживания партитуры}
состоит в построении \textit{каузального} отображения
\begin{equation}
\varphi : \{1, \dots, M\} \longrightarrow \{1, \dots, N\},
\qquad j \mapsto \varphi(j),
\eqlab{alignment-formal}
\end{equation}
ставящего в соответствие каждому событию исполнения с номером~$j$
позицию $\varphi(j)$ в партитуре, причём оценка $\hat\varphi(j)$
вычисляется по информации только из прошлого: $q_1, \tau_1,
\dots, q_j, \tau_j$. Поскольку в большинстве музыкальных
контекстов исполнитель движется по партитуре вперёд~--- с
возможными редкими пропусками, повторами и
ошибками~\cite{Nakamura2015}, но без систематического возврата
назад~--- отображение~$\varphi$ обычно требуют удовлетворяющим
условию монотонности
\begin{equation}
j_1 < j_2 \Longrightarrow \varphi(j_1) \le \varphi(j_2).
\eqlab{monotonicity-formal}
\end{equation}

Для практического применения функция выравнивания дополняется
оценкой текущего темпа исполнения $\hat v(t) > 0$, выражающей
скорость движения исполнителя по партитуре в beat/секунду.
Темп используется блоком воспроизведения сопровождения для
синхронного синтеза оркестровой партии. В наиболее общей
постановке~\cite{Cont2010,Raphael2010} обе величины оцениваются
совместно, а не последовательно: апостериорное распределение по
позиции зависит от текущей оценки темпа, а оценка темпа~--- от
истории согласованных событий.

\paragraph*{Дополнительные требования.}
\textit{Каузальность:} оценка $\hat\varphi(j)$ должна быть
доступна не позже момента $\tau_j + L$, где $L$~--- допустимая
задержка системы.
\textit{Робастность:} ошибки исполнителя (пропуски нот,
фальшивые ноты, орнаменты), не предусмотренные партитурой,
не должны приводить к потере синхронизации.
\textit{Адаптивность:} оценка темпа должна реагировать на
выразительные изменения исполнителя~--- ритардандо, ачелерандо,
rubato~--- не теряя устойчивости.
